\documentclass[conference]{IEEEtran}
\usepackage{amsmath}
\usepackage{booktabs}
\usepackage{array}
\usepackage{graphicx}
\usepackage{hyperref}
\usepackage{xcolor}
\usepackage{multirow}

\newcommand{\ci}[2]{[#1, #2]}
\sloppy
\emergencystretch=3em
\clubpenalty=10000
\widowpenalty=10000

\begin{document}

\title{Comparing Rule-Based, Game-Theoretic, and Reinforcement-Learned\\
Response Selection for Residential IoT Defense:\\
A Controlled Mininet Evaluation with Disclosed Self-Correction}

\author{
\IEEEauthorblockN{Author Name\IEEEauthorrefmark{1}}
\IEEEauthorblockA{\IEEEauthorrefmark{1}Affiliation \\
Email: author@example.com}
}

\maketitle

\begin{abstract}
Automated response selection for residential IoT networks is typically evaluated either as a
detection problem (is the traffic malicious?) or through anecdotal demonstration, rarely as a
controlled comparison of \emph{how the response itself is chosen}. We present a registry-driven
defense system for a Mininet-simulated residential IoT network that compares three independent
response-selection policies -- a rule-based classifier, a Stackelberg game-theoretic solver, and
a PPO-trained reinforcement-learning agent -- against two naive baselines and an external
signature-based IDS (Suricata), on real captured traffic across 17 threat conditions (16 attack
types plus benign). Our evaluation harness runs 255 real Mininet trials per reported round
(1{,}275 recorded outcomes), reporting Wilson-score confidence intervals rather than point
estimates, and includes ablations on the Stackelberg payoff table, PPO reward shaping,
training-seed robustness, and decision-boundary policy divergence, backed by a formal paired
significance test (McNemar's exact test) rather than agreement-rate percentages alone. What
distinguishes this evaluation from a typical system paper is not its headline accuracy but the
process that produced it. Re-running the harness at three times an initial sample size first
revealed that several headline numbers had been optimistic at low $N$ (aggregate detection
accuracy fell from $90.0\%$ to $73.7\%$); rather than accept that as a ceiling, we investigated
each underperforming condition as a live, reproducible bug, found and fixed seven real defects
(a capture-race condition, a traffic-pacing error, a silently-wrong protocol default, two
Suricata-ruleset gaps, a reward-function asymmetry, and a packet-buffering bug), and re-ran the
full harness under the fixed code: aggregate detection accuracy reached $99.6\%$
\ci{97.8}{99.9}, with $16$ of $17$ conditions scoring a clean $100\%$. One further attempted fix
-- quadrupling PPO's training budget to eliminate a residual seed-dependent failure -- was tried,
measured on a full ten-seed sweep, found not to generalize from the single seed it was validated
on, and reverted; we report this as a documented dead end, not a success. We conclude that the
three response-selection policies are behaviorally indistinguishable on in-distribution traffic
(McNemar's exact test, $p=1.0$) but diverge significantly near decision boundaries
($p<1.2\times10^{-159}$), that our system now materially outperforms naive baselines by a wide,
formally-tested margin, and that every result in this paper is scoped to a single controlled VM
and Mininet topology -- external validity remains an explicitly open problem, and internal
validity was earned through disclosed, repeated self-correction rather than assumed.
\end{abstract}

\begin{IEEEkeywords}
IoT security, intrusion response, Stackelberg games, reinforcement learning, Mininet, reproducibility
\end{IEEEkeywords}

\section{Introduction}

Consumer and residential IoT devices are widely deployed with weak or absent built-in security,
and network-level defense -- detecting an attack and then \emph{choosing what to do about it} --
is increasingly delegated to automated systems. The detection side of this problem is well
studied: signature-based tools such as Snort and Suricata, and a large body of machine-learning
intrusion-detection literature, address the question of whether observed traffic is malicious.
The \emph{response-selection} side -- given a detected threat, which of several possible actions
(alert, isolate, deceive, rate-limit, quarantine, \ldots) should actually be taken -- receives
comparatively little controlled, comparative evaluation. Individual mechanisms (game-theoretic
defense, reinforcement-learned defense, moving-target defense) are each well represented in the
literature, but rarely compared against each other, against simple baselines, and against an
external detector, on the same captured traffic, in the same controlled environment.

This paper makes four contributions:

\begin{enumerate}
\item A registry-driven defense architecture for a simulated residential IoT network (Section~III)
  in which a single declarative record per attack type drives traffic generation, detection,
  Stackelberg payoff lookup, PPO observation encoding, and dataset generation, so that adding a
  new attack type or defense action does not require touching each subsystem by hand.
\item A repeatable benchmark harness (Section~IV) that runs three independent response-selection
  policies -- rule-based, Stackelberg game-theoretic, and PPO reinforcement-learned -- alongside
  two naive baselines and an external, real signature-based IDS, against the \emph{same} captured
  Mininet traffic per trial, with every chosen action actually executed and independently verified
  rather than assumed successful.
\item A set of ablations and robustness checks (Section~V) largely absent from comparable systems
  work: Wilson-score confidence intervals on every reported rate, a payoff-table sensitivity
  ablation, a reward-shaping ablation, a ten-seed PPO training-robustness sweep, a direct test of
  policy divergence at synthetic decision boundaries, and a formal paired significance test
  (McNemar's exact test) confirming that divergence rather than resting on agreement-rate
  percentages alone.
\item A transparent account of three self-corrections (Section~\ref{sec:corrections}): a threefold
  sample-size increase that first revealed several headline numbers had been optimistic, followed
  by a systematic bug-hunt that found and fixed six real defects and recovered a $99.6\%$ aggregate
  accuracy; an adaptive-attacker scale-up that discovered a real bug in our own evaluation code and,
  once fixed, failed to reproduce an earlier finding; and a training-budget fix that resolved one
  isolated failure but did not generalize on full re-test, and was reverted. All three are reported
  as corrections, including the one that failed after already being committed, rather than
  suppressed or quietly kept.
\end{enumerate}

We deliberately scope every claim in this paper to the environment it was measured in: a single
Mininet topology on one resource-constrained virtual machine (Section~VII). We do not claim
external validity, and we treat the absence of that claim as a feature of a rigorous internal
evaluation, not an oversight to be hidden.

\section{Related Work}

No published study was found comparing rule-based, Stackelberg game-theoretic, and
reinforcement-learned \emph{response selection} for IoT defense on a shared dataset; the closest
work in the literature compares detection methods alone, or evaluates a single game-theoretic or
learned defense mechanism on its own metrics and environment. We situate our design choices
against this literature as context, not as a claimed head-to-head benchmark.

\textbf{Signature-based detection.} An experimental comparison of Suricata and Snort found
Suricata generally ahead on detection accuracy, scalability, and resource efficiency (e.g.\ $100\%$
vs.\ $85.7\%$/$66.7\%$ on two DNS-tunneling variants) \cite{nps_suricata_snort}, motivating our
choice of Suricata as the external detection baseline (Section~IV-C). Prior work evaluating
signature-based tools against IoT botnet datasets (ISOT, IoT-23, Bot-IoT) documents the same
structural limitation we independently observe: signature/threshold detectors miss traffic that
does not match a known pattern \cite{iot_botnet_detection}.

\textbf{Game-theoretic and deception-based defense.} A stochastic-game honeypot moving-target-defense
study reports substantial real engagement gains from deception over a static defense (attacker
packet interactions up $920$ vs.\ $42$; scanner engagement extended $14.6\times$)
\cite{mtd_honeypot_stochastic_game} -- a different environment and metric than our controlled,
single-attacker Mininet trials, but directionally consistent with our own DECOY action's measured
value (Section~V-B). A 2025 paper models collaborative IoT packet-sampling against DDoS as a
Stackelberg game and derives a sampling-rate lower bound that deters an attacker
\cite{stackelberg_iot_ddos_sampling} -- the same game-theoretic framing we use, but for a
continuous sampling-rate decision rather than a choice among discrete response actions, and
evaluated analytically rather than against captured traffic.

\textbf{Reinforcement learning for defense.} Recent surveys document deep RL as an active,
promising direction for IoT intrusion detection while flagging reproducibility and dataset-realism
concerns across the field \cite{drl_iot_survey} -- concerns we address directly by training and
evaluating our PPO policy against real, verified Mininet outcomes and by reporting a
seed-robustness sweep (Section~\ref{sec:seedrobustness}) rather than a single trained checkpoint. A
2025 paper evaluates an RL-tuned moving-target-defense mutation strategy for edge IoT against DDoS
specifically on attack success rate, defense latency, and CPU overhead \cite{rl_mtd_edge_iot} --
the same category of cost-conscious framing our own cost/overhead analysis uses
(Section~\ref{sec:cost}), for mutating network
configuration rather than selecting among discrete response actions.

None of the above papers shares our dataset, environment, or exact decision space, which is
precisely why our results are reported as an internally-valid comparison between the arms tested
here, not translated into a claimed standing against any of them.

\section{System Architecture}

\subsection{Registry-driven attack and response model}

The system is organized around a single declarative registry
(\texttt{AttackScenario}) with one record per attack type, holding its traffic generator,
rule-based detector, registered ``preferred'' response action, and the confidence/score thresholds
that action requires. Every dependent subsystem -- the rule-based decision policy, the Stackelberg
payoff lookup, PPO's one-hot scenario encoding, and dataset generation for the machine-learning
baseline -- derives its behavior from this registry rather than encoding per-attack logic
separately in each place. This was a deliberate response to an early failure mode in the project's
own development: adding one new attack type by hand across six independently-maintained files
produced four distinct, reproduced bugs (a malformed traffic-generation script that hung
indefinitely, an insufficient process-termination grace window, a privileged-port binding failure,
and output-buffering interference with the network simulator's own control channel). The registry
refactor eliminates that failure class by construction for every attack added afterward.

\subsection{Threat model and attack catalog}

Table~\ref{tab:attacks} summarizes the 16 registered attack types plus the benign (\texttt{normal})
condition, each with a distinct traffic signature and a registered preferred response. Signatures
span packet rate, port diversity, payload size, traffic direction (several attacks, e.g.\
\texttt{data\_exfiltration} and \texttt{dns\_tunneling\_exfiltration}, reverse the usual
device-is-target direction), and, for \texttt{c2\_beaconing}, inter-packet timing regularity --
the only signature in the registry that is not rate-, volume-, or size-based. Every detector's
numeric threshold window was placed in a verified, unclaimed gap relative to every other
detector's window, confirmed both by a parametrized test sweeping every attack's synthetic
signature through the complete detector set and by repeated live Mininet runs under real timing
variance.

\begin{table}[t]
\centering
\caption{Registered threat conditions and preferred responses}
\label{tab:attacks}
\scriptsize
\setlength{\tabcolsep}{3pt}
\begin{tabular}{@{}p{0.34\linewidth}p{0.30\linewidth}p{0.24\linewidth}@{}}
\toprule
\textbf{Condition} & \textbf{Signature axis} & \textbf{Preferred} \\
\midrule
normal & -- (benign baseline) & ALLOW \\
reconnaissance port scan & many ports, moderate rate & DECOY \\
DoS flood & one port, very high rate & ISOLATE \\
brute force & one port, sustained attempts & BLOCK\allowbreak\_SOURCE \\
data exfiltration & reversed direction, large payload & ISOLATE \\
exploit payload injection & oversized single payload & DECOY \\
TCP SYN flood & refused connects, one port & ISOLATE \\
ICMP ping flood & high-volume padded ICMP & THROTTLE \\
slow-loris exhaustion & many held-open connections & THROTTLE \\
DNS amplification & oversized inbound UDP & BANDWIDTH\allowbreak\_CAP \\
DNS tunneling exfiltration & reversed, many small UDP & DECOY \\
MQTT message flood & many completed TCP conns. & THROTTLE \\
firmware tampering & reversed, periodic oversized push & QUARANTINE \\
buffer overflow probe & sustained oversized payload, 1 conn. & RESET\allowbreak\_SESSIONS \\
credential replay & repeated uniform small UDP & BLOCK\allowbreak\_SOURCE \\
rogue config beacon & reversed, low-confidence pattern & FORENSIC\allowbreak\_CAPTURE \\
C2 beaconing & periodic inter-arrival timing & ISOLATE \\
\bottomrule
\end{tabular}
\end{table}

\subsection{Defense actions}

Ten actions are available to all three policies, all real, verified mechanisms executed inside
the Mininet network rather than simulated effects: \texttt{ALLOW} and \texttt{ALERT} (no
enforcement); \texttt{ISOLATE} (interface down, verified by polling link state);
\texttt{DECOY} (a real TCP banner service with NAT redirection, verified by a real redirected
connection); \texttt{THROTTLE} (a protocol-aware \texttt{iptables hashlimit} rate cap, verified
against real traffic -- an unthrottled attacker completed $56/56$ connection attempts in $3\,$s,
the same traffic under a $2/$s cap completed $7$); \texttt{BLOCK\_SOURCE} (a source-scoped
\texttt{iptables} DROP rule, verified by polling the installed rule); \texttt{QUARANTINE}
(default-deny with an allowlist for known-legitimate hosts); \texttt{RESET\_SESSIONS} (real kernel
socket destruction via \texttt{ss~-K}); \texttt{FORENSIC\_CAPTURE} (preserves the detection
pipeline's own already-captured evidence); and \texttt{BANDWIDTH\_CAP} (a \texttt{tc tbf} egress
cap on the \emph{attacker's} interface, verified at the receiver: $300/300$ packets delivered
uncapped vs.\ $20/300$ under the cap).

\subsection{Response-selection policies}

\textbf{Rule-based.} A deterministic mapping from the classified attack type and its
score/confidence to the registry's preferred action, with dedicated handling for the benign and
low-confidence cases.

\textbf{Stackelberg.} A leader-follower game solved per detected threat type, with a full,
independently-reasoned payoff table (attacker and defender utility for every action, for every
threat type, in \texttt{config/policies.yaml}) rather than a derived or auto-generated one. The
solver's own selection was verified to agree with the registry's declared preferred action for
all 17 conditions.

\textbf{PPO.} A proximal-policy-optimization agent (Stable-Baselines3) over a ten-action discrete
space matching \texttt{DefenseAction}'s declaration order, trained against a deterministic
normalized security-context observation. The deployed recipe is \texttt{net\_arch=[64,64]},
\texttt{ent\_coef=0.01}, \texttt{n\_steps=batch\_size=170}, learning rate $0.001$, $25{,}500$
timesteps. An initial \texttt{net\_arch=[32,32]} network reliably mis-routed the benign scenario
to DECOY regardless of training length; doubling network capacity resolved this at the project's
development seed. Section~\ref{sec:seedrobustness} reports a formal, multi-seed robustness check of this recipe rather
than relying on that single seed's success.

\section{Evaluation Methodology}

\subsection{Harness design}

Our benchmark harness runs $N$ real Mininet trials for every registered condition. Each trial (i)
generates real traffic for that condition (the same generators the live system uses) and captures
it, preserving the raw packet capture for later offline analysis; (ii) classifies the capture with
the attack-type-agnostic rule-based detector; (iii) evaluates five arms against the
\emph{identical} resulting security context -- the rule-based, Stackelberg, and PPO policies, plus
\texttt{NaiveBlockAllBaseline} (isolate any detected threat, the common ``one fixed response to
every alert'' design) and \texttt{AlwaysAllowBaseline} (never intervene, the floor); and (iv)
executes and independently verifies each \emph{distinct} action chosen across the five arms
exactly once, for real. One trial produces every arm's result; no arm requires a separate live
run.

\subsection{Metrics}

\emph{Detection accuracy}: did the shared detector correctly classify the captured traffic (one
result per trial, shared across arms, reported as a sanity check rather than a point of
comparison). \emph{Matches preferred action}: of the attack trials, how often an arm's chosen
action was the registry's declared preferred action. \emph{Verified response rate}: of all trials,
how often an arm both chose the preferred action \emph{and} that action's real effect was
independently confirmed -- our headline metric, zero whenever the wrong action was chosen, only
crediting a checked success. Every reported rate carries a $95\%$ Wilson-score confidence
interval, computed directly rather than via the normal approximation (which misbehaves at small
$n$ or extreme proportions).

\subsection{External detection baseline}

A second, independent arm runs offline Suricata analysis of the same captured packet captures,
against two rulesets: ET-Open (real-world community threat signatures) and a small lab-tailored
ruleset encoding the same shape knowledge our own rule-based detectors use, so Suricata is judged
on a fair signature-matching version of that knowledge too, not only on a strawman of unrelated
real-world malware signatures. Suricata is compared on detection axis only (accuracy,
true-/false-positive rate against ground truth); it identifies traffic but does not select or
execute a containment action.

\subsection{Scale}

This paper reports the harness at $\text{trials\_per\_condition}=15$ across the full 17-condition
registry: $255$ real Mininet trials, $1{,}275$ recorded outcomes, on 2026-09-25 ($\approx$2h21m
wall-clock, run after all seven fixes in Section~\ref{sec:corrections} landed). Two earlier runs at
this same scale -- one at $N=5$, one at $N=15$ but against the pre-fix code
(Section~\ref{sec:sizecorrection}) -- are discussed explicitly because their headline numbers did
not hold, first at the larger sample and then again once the underlying code was corrected.

\section{Results}

\subsection{Main comparison}

Table~\ref{tab:main} reports the five-arm comparison \emph{after} all seven bugs in
Section~\ref{sec:corrections} were fixed and the harness was re-run in full. The three real
policies are statistically indistinguishable from each other on this dataset (confirmed formally
in Section~\ref{sec:mcnemar}) and substantially outperform both baselines.

\begin{table}[t]
\centering
\caption{Five-arm comparison, $N=15$ trials/condition ($n=255$ trials), post-fix}
\label{tab:main}
\scriptsize
\setlength{\tabcolsep}{3pt}
\begin{tabular}{@{}lccc@{}}
\toprule
\textbf{Arm} & \textbf{Detection acc.} & \textbf{Verified resp.} & \textbf{Matches pref.} \\
\midrule
Rule-based (ours) & 99.6\ci{97.8}{99.9} & 99.6\ci{97.8}{99.9} & 99.6\ci{97.7}{99.9} \\
Stackelberg (ours) & 99.6\ci{97.8}{99.9} & 99.6\ci{97.8}{99.9} & 99.6\ci{97.7}{99.9} \\
PPO (ours) & 99.6\ci{97.8}{99.9} & 99.6\ci{97.8}{99.9} & 99.6\ci{97.7}{99.9} \\
NaiveBlockAll & 99.6\ci{97.8}{99.9} & 29.4\ci{24.2}{35.3} & 25.0\ci{19.9}{30.8} \\
AlwaysAllow & 99.6\ci{97.8}{99.9} & 5.9\ci{3.6}{9.5} & 0.0\ci{0.0}{1.6} \\
\bottomrule
\end{tabular}
\end{table}

All values are percentages with $95\%$ Wilson CIs. Sixteen of the seventeen conditions score a
clean $15/15$; the one exception, \texttt{credential\_replay} ($93.3\%$), has a completely empty,
unreadable pcap on its one miss -- a pre-existing, disclosed capture-flakiness class distinct
from the capture-race bug fixed in Section~\ref{sec:corrections}. Our system's $99.6\%$
verified-response rate is $3.4\times$ \texttt{NaiveBlockAllBaseline}'s $29.4\%$ and $16.9\times$
\texttt{AlwaysAllowBaseline}'s $5.9\%$ -- a wider margin than any earlier-scale estimate, because
our own numbers improved dramatically while the baselines, unaffected by any of the seven fixes,
stayed roughly where they were. Mean detection latency ($31{,}287\,$ms) is dominated by real
packet capture, not decision-making (Section~\ref{sec:cost}).

\subsection{External detection comparison}

Table~\ref{tab:suricata} compares our detector against Suricata under both rulesets, after fixing
\texttt{dos\_flood}'s non-firing rule and writing real signatures for the eleven attacks the lab
ruleset previously lacked (Section~\ref{sec:corrections}).

\begin{table}[t]
\centering
\caption{External detection comparison, offline analysis of 255 pcaps, post-fix}
\label{tab:suricata}
\scriptsize
\setlength{\tabcolsep}{3pt}
\begin{tabular}{@{}lccc@{}}
\toprule
\textbf{Detector} & \textbf{Accuracy} & \textbf{TPR} & \textbf{FPR} \\
\midrule
Suricata + ET-Open & 43.9\% & 41.3\% & 13.3\% \\
Suricata + lab rules & 94.1\% & 93.8\% & 0.0\% \\
Ours (rule-based) & 99.6\ci{97.8}{99.9} & -- & 0.0\% \\
\bottomrule
\end{tabular}
\end{table}

The lab-tailored ruleset now performs close to our own detector -- real, independent, external
validation of the fixed signatures. ET-Open's own numbers moved too (previously $67.8\%$
accuracy) despite nothing in this pass touching its ruleset: a real, disclosed side effect of
several traffic generators changing shape this session (Section~\ref{sec:corrections}), not a new
defect -- ET-Open's community signatures were never tuned for this lab's synthetic traffic in the
first place.

\subsection{Supplementary ML confirmation layer}

Separately from the three response-selection policies, a Random Forest classifier serves as a
confirmation step for the one signature with the fuzziest rule-based boundary
(\texttt{reconnaissance\_port\_scan}); it is not itself one of the compared policies. Its own
training dataset was, until this pass, limited to the original 6-class registry, predating the
expansion to 17 conditions. We generated a fresh, real 17-class dataset ($255$ Mininet runs, $0$
failures, $260$ rows balanced at $15$/class) and retrained it: on a held-out test split ($40$
rows), the Random Forest reaches $97.5\%$ accuracy (precision $91.2\%$, recall $92.2\%$) against a
same-split rule-based baseline of $90.0\%$ -- a real, measured $7.5$-point gap on the full registry,
superseding an earlier $6$-class figure that predated the ten-attack expansion. Two classes showed
$0$ precision/recall on this specific split at only $2$--$3$ held-out rows each; reported as
small-sample noise, not a claimed per-class weakness.

\subsection{Cost of comparing three policies}
\label{sec:cost}

Timing $300$ real, direct calls per policy (isolated from Mininet and packet capture) gives mean
decision latencies of $0.054\,$ms (rule-based), $0.464\,$ms (Stackelberg), and $1.136\,$ms (PPO).
Running all three sequentially, as our comparison architecture does, costs $1.654\,$ms total; the
marginal overhead of comparing three policies versus deploying only the slowest one is
$0.518\,$ms -- approximately $0.0017\%$ of the harness's own $31.3\,$s mean real detection latency.
Comparing multiple policies is not this system's bottleneck.

\subsection{PPO training seed-robustness}
\label{sec:seedrobustness}

A single trained checkpoint is a weak basis for a robustness claim. We trained 10 independent
models at the deployed recipe, varying only the random seed
($1,2,3,7,13,17,42,55,88,99$), and checked full convergence to every registered preferred action.
An initial sweep found $8/10$ fully converge (Wilson $95\%$ CI \ci{49.0}{94.3}), both failures the
same class of mistake -- a false-positive-style error on \emph{benign} traffic (one seed chose
THROTTLE, one chose DECOY, instead of ALLOW). Tracing this to a real reward-function asymmetry
(only ISOLATE-on-normal carried an extra disruption penalty, giving the optimizer less pressure to
avoid the other two actions specifically) and fixing it raised the deployed recipe's own rate to
$\mathbf{9/10}$ (Wilson $95\%$ CI \ci{59.6}{98.2}) -- the figure Table~\ref{tab:main} was trained
under. A further attempt to eliminate the one remaining failure did not hold up under full
re-testing; Section~\ref{sec:budgetcorrection} reports it as a documented dead end. The wide CI is
reported honestly rather than rounded into overconfidence: this supports ``most seeds converge
fully; one specific instability mode recurs in roughly one run in ten,'' not a tighter claim.

\subsection{Ablations}

\textbf{Stackelberg payoff sensitivity.} We perturbed every value in the deployed payoff table by
random noise proportional to its own magnitude ($\pm5\%$ to $\pm50\%$, 30 trials/level, 2{,}550
solves total) and measured the fraction of solves still matching the registry's declared action.
Table~\ref{tab:stackelberg} shows a monotonic, graceful degradation: small realistic tuning error
barely moves the outcome (the qualitative payoff \emph{structure} is doing most of the work), while
at $\pm50\%$ noise over a third of solves diverge, confirming the specific values still carry real
information.

\begin{table}[t]
\centering
\caption{Stackelberg payoff-table sensitivity (2{,}550 solves)}
\label{tab:stackelberg}
\small
\begin{tabular}{@{}lc@{}}
\toprule
\textbf{Noise level} & \textbf{Match rate (95\% CI)} \\
\midrule
$\pm5\%$  & 96.1\ci{94.0}{97.4} \\
$\pm10\%$ & 88.0\ci{85.1}{90.5} \\
$\pm20\%$ & 78.8\ci{75.3}{82.0} \\
$\pm35\%$ & 64.5\ci{60.5}{68.4} \\
$\pm50\%$ & 61.0\ci{56.8}{65.0} \\
\bottomrule
\end{tabular}
\end{table}

\textbf{PPO reward shaping.} We trained the real environment under two reward configurations,
holding the recipe and a five-seed sweep fixed: the deployed \emph{shaped} reward (different
magnitudes for different outcome severities, e.g.\ letting an attack through unopposed is
penalized twice as hard as choosing a merely suboptimal defensive response) against a \emph{flat}
reward collapsing every correct outcome to the same reward and every incorrect outcome to the same
penalty. Both configurations reached full convergence on $5/5$ seeds (Wilson CI \ci{56.6}{100.0}
for both). We report this honestly as inconclusive at this sample size, not as evidence that
shaping is unnecessary: both intervals are wide and fully overlapping, and the more likely places
shaping would matter -- sample efficiency, or a longer-horizon task than our single-decision
environment -- were not measured here.

\subsection{Decision divergence near boundaries}
\label{sec:divergence}

Because real captured traffic in our lab classifies far from any threshold every time, the three
policies choosing identical actions on Table~\ref{tab:main} is closer to a consistency check than
a comparison. We separately tested all three policy classes directly against synthetic security
contexts: (i) threat-score/confidence perturbed around each attack's own registered decision
threshold ($\pm0.03$ to $\pm0.15$, $784$ contexts across the full registry), and (ii) entirely
unregistered ``novel'' attack types no policy was tuned for ($12$ contexts). Near decision
boundaries, rule-based diverges from both Stackelberg and PPO on roughly two-thirds of ambiguous
contexts (all-three agreement $32.5\%$), while Stackelberg and PPO track each other closely
($100\%$ mutual agreement) even off the training distribution -- a positive finding that the
learned policy generalizes toward the same strategic reasoning the game solver encodes explicitly.
On entirely unregistered attack types, rule-based and Stackelberg both degrade to a safe ALERT
default every time; \textbf{PPO instead selects ISOLATE regardless of the actual threat score or
confidence supplied} -- a genuine, disclosed limitation: PPO has no mechanism for recognizing
``I have never seen anything like this,'' unlike the other two policies' explicit fallback.

\subsection{Formal significance testing: McNemar's exact test}
\label{sec:mcnemar}

Every comparison above reports agreement or verified-response \emph{percentages} -- real, but not
a formal test of whether an observed difference could plausibly be chance. Because every arm is
scored against the identical trial or synthetic context, this is a \emph{paired} design, and
McNemar's test is the appropriate significance test for it: it conditions on the trials where two
policies actually disagree, since a trial both get right (or both get wrong) carries no
information about which is \emph{better}. We implement the exact binomial form rather than the
more commonly cited chi-square approximation, which is unreliable at the small discordant-pair
counts a paired, mostly-agreeing comparison like this one produces, and verify our implementation
against a hand-computed reference before applying it to real data.

On the $255$ real harness trials (Table~\ref{tab:main}), all three pairwise comparisons show
\emph{zero} discordant pairs ($p=1.0$) -- the rigorous version of a claim this paper has otherwise
only supported with matching percentages: on real captured traffic, there is no statistical
evidence whatsoever of a performance difference between any pair of policies, because the three
never once chose differently across these 255 trials. On the $784$-context synthetic
decision-boundary dataset (Section~\ref{sec:divergence}), where real disagreement exists, the same
test tells a sharply different story: both Stackelberg and PPO significantly outperform rule-based
($529/784$ discordant pairs, $100\%$ in their favor, exact $p<1.2\times10^{-159}$), while remaining
statistically indistinguishable from each other ($p=1.0$, zero discordant pairs) -- formal
confirmation, not just an observed agreement rate, that the learned policy converges toward the
same strategic reasoning the game-theoretic solver encodes explicitly.

\subsection{Detector robustness (evasion margin)}
\label{sec:evasion}

Perturbing every numeric feature in each attack's own signature by $-30\%$ to $+30\%$ and
re-classifying through the complete detector set measures brittleness versus real margin. Results
are uneven, not uniformly reassuring: robust classification rates range from $28.6\%$
(\texttt{dns\_amplification}, \texttt{firmware\_tampering}, \texttt{buffer\_overflow} -- a crowded
region of the numeric feature space) to $100\%$ (\texttt{dos}, \texttt{brute\_force},
\texttt{exfiltration}). We are explicit that this measures \emph{measurement-noise tolerance}, not
\emph{adversarial evasion resistance}: the perturbation is symmetric, random, and non-adaptive, and
does not search for the nearest evasive point the way a real, threshold-aware adversary would
(Section~VII).

\section{Discussion: Self-Correction as Methodology}
\label{sec:corrections}

We report three findings in this section specifically because our first attempt at strengthening
each one did not simply add more of the same evidence -- it changed what we could honestly claim,
in one case for the better once the underlying cause was fixed, and in another not at all.

\subsection{Sample size revealed optimism -- and investigation fixed it}
\label{sec:sizecorrection}

An earlier harness run at $N=5$ trials/condition across 16 conditions reported aggregate detection
accuracy of $90.0\%$ and verified-response rate of $85.0\%$. Re-running at $N=15$ across the full
17-condition registry revised both down substantially, to $73.7\%$ and $68.2\%$. Several individual
conditions that read as a clean $100\%$ at $N=5$ settled meaningfully lower at $N=15$ -- exactly the
outcome a single unlucky trial swinging a small-sample percentage by a full $20$ points predicts.
Two conditions illustrated this most sharply: \texttt{tcp\_syn\_flood} fell to $20.0\%$ detection
accuracy, and \texttt{icmp\_ping\_flood} to a reproducible $0\%$ verified-response rate (CI
\ci{0.0}{20.4}, now excluding anything above roughly one-fifth).

Rather than accept $73.7\%$ as this system's ceiling, we treated every underperforming condition as
a live, reproducible bug and investigated each with a direct repro rather than a post-hoc
explanation. Six real, distinct defects were found and fixed: (1) a capture-race condition in which
non-unique capture filenames let a not-fully-terminated \texttt{tcpdump} process from one round
silently corrupt the next round's capture on any persistent-network module; (2) \texttt{tcp\_syn\_flood}'s
traffic generator targeting the wrong packets-per-second rate, compounded by a fixed
post-\texttt{connect()} sleep that let real system-call latency dominate the achieved rate --
fixed with a deadline-based scheduler, verified live at $17.58$--$17.61$ packets/sec across six
consecutive trials, all six correctly classified (was $3/15$); (3) \texttt{THROTTLE} silently
defaulting to a TCP-only firewall rule for every attack because the caller populated an empty
context, never wiring through the real detected protocol -- the exact and sole cause of
\texttt{icmp\_ping\_flood}'s $0\%$ figure, confirmed end-to-end before and after the fix; (4) a
Suricata rule type (\texttt{threshold}) that never once fires for a UDP flow in our environment,
confirmed by testing it in isolation against four real pcaps with a count far below what the
traffic contained; (5) our lab-tailored Suricata ruleset covering only $5$ of $17$ attacks, plus a
second, compounding discovery that a partial fix already existed in this project's own history but
had never been deployed to the path Suricata actually reads from; and (6) a capture-buffering bug
in which \texttt{tcpdump}'s default behavior only flushes to disk once its buffer fills, losing the
final packet of any sufficiently sparse capture if a forced termination is ever needed -- fixed with
tcpdump's own \texttt{-U} (packet-buffered) flag.

Every fix was verified with a live, direct repro before being accepted, and the full harness was
then re-run under the fixed code (Table~\ref{tab:main}): aggregate detection accuracy reached
$99.6\%$, with $16$ of $17$ conditions scoring a clean $100\%$. We report this as the central
methodological result of this paper, not a footnote to it: a $16$-point accuracy revision is not,
on its own, evidence of a well-run study -- treating that revision as a set of falsifiable
hypotheses, rather than an accepted ceiling, is.

\subsection{An adaptive-attacker claim that did not survive scaling}
\label{sec:adaptivecorrection}

A separate module tests whether identity-based containment (\texttt{BLOCK\_SOURCE}) actually
contains a persistent attacker across repeated attempts, in two modes: \emph{fixed} (the same
source every round, a control) and \emph{rotate} (a fresh, not-yet-blocked source every round). An
initial five-round pilot suggested a real finding: a second, fresh spoofed source (\texttt{.121})
reached the sensor and was independently detected and blocked on its first attempt, consistent
with \texttt{BLOCK\_SOURCE}'s identity-scoped mechanism (confirmed directly from its
implementation) offering no protection against a rotating attacker.

We scaled this to 30 rounds/mode to obtain more than two data points. The scaled run did not
produce the expected volume of clean evidence -- instead, reading the raw per-round output
revealed a real bug in our own module: rotate-mode's source selection always chose the first
not-yet-blocked address in a fixed pool. Because the second address's capture came back empty
(ambiguous between ``blocked'' and ``not observed,'' a distinction our own retry logic is designed
to disclose rather than conflate), it stayed first-in-line and was re-selected for all 28 remaining
rounds instead of the campaign rotating through the rest of the pool. We fixed the selection logic,
added a regression test, expanded the source pool from 6 to 21 addresses, and re-ran (25
rounds/mode, genuinely cycling through 20 distinct spoofed addresses this time, confirmed from the
log). \textbf{Every one of the 20 distinct spoofed addresses produced zero captured packets, on
both the initial attempt and the built-in retry (40 real zero-packet observations), in two
independent bug-fixed runs.} Only the physical attacker's own real, unspoofed address ever
produced capturable traffic.

The honest conclusion this now supports is weaker than the original pilot suggested: the one
successful spoofed capture that grounded our original claim could not be reproduced across 40
further attempts. We report this as \emph{one unreplicated pilot observation, not confirmed by two
larger, bug-fixed follow-up attempts} -- not as strengthened evidence for source-rotation evasion.
The architectural argument (\texttt{BLOCK\_SOURCE} is scoped to the packet's claimed source
address, not the physical sender) still stands independent of this measurement, but we no longer
claim repeated Mininet evidence that a rotating attacker gets through. Root-causing why spoofed
traffic is not reliably captured in this environment (host-level reverse-path filtering was
checked and does not explain it; a per-namespace or virtual-switch-level cause was not further
isolated) is left as future work.

\subsection{A training-budget fix that looked real and did not generalize}
\label{sec:budgetcorrection}

Section~\ref{sec:seedrobustness} reports a reward-function fix that raised PPO's seed-convergence
rate from $8/10$ to $9/10$, leaving one seed (17) still failing on the identical mistake. Isolating
that one seed and retraining it at $2\times$ and $4\times$ the deployed timestep budget showed the
mismatch move to a different scenario at $2\times$ and disappear entirely at $4\times$ -- a result
that, tested on that one seed alone, looked like a genuine, targeted fix. Generalizing from it was
premature: a full ten-seed sweep at the $4\times$ budget came back to $8/10$, not an improvement.
Seed 17 did converge, exactly as the isolated test predicted, but seeds 3 and 13 -- both fully
converged at the original budget -- newly failed on two scenarios that had never been a problem for
any seed before. Quadrupling the training budget did not fix the aggregate seed-fragility; it moved
which seeds and scenarios were most sensitive to initialization, at four times the real training
cost. We reverted the change. The deployed recipe and the $9/10$ figure in
Section~\ref{sec:seedrobustness} are the reward fix alone, at the original training budget -- the
best real recipe this investigation found, not the most recent one tried.

\section{Limitations}

\textbf{Single environment.} Every result in this paper comes from one Mininet topology on one
resource-constrained ($2$-CPU, $\approx2.9\,$GB) virtual machine. No claim here should be read as
generalizing to arbitrary IoT hardware, network topologies, or attacker populations; we have not
tested beyond this single environment.

\textbf{No adaptive, threshold-aware attacker.} Our robustness measurements (Section~\ref{sec:evasion}) use
symmetric, non-adaptive perturbation. A real adversary that probes and adapts to our specific
detection thresholds was not modeled.

\textbf{Shared detection across internal arms.} The three policies and two baselines receive
identical detection results per trial; our comparison measures response selection and execution,
not detection accuracy, which is why the Suricata arm exists as an independent check.

\textbf{Adaptive-attacker evidence is explicitly weak}, as detailed in
Section~\ref{sec:adaptivecorrection} -- an unreplicated pilot observation, not a demonstrated finding.

\textbf{One condition retains a real, disclosed residual gap.} \texttt{credential\_replay} scores
$93.3\%$ ($14/15$); its one miss has a completely empty, unreadable packet capture -- a general
Mininet/\texttt{tcpdump} capture-flakiness class distinct from the capture-race bug fixed in
Section~\ref{sec:sizecorrection}, and not further resolved in this pass.

\section{Ethics and Responsible Disclosure}

Every attack technique implemented in this system targets only hosts created inside an isolated
Mininet virtual network on a single VM under our own control, with no path to any real external
host or network. A systematic self-review of this codebase found and fixed 21 issues, two of which
would be real security weaknesses in a production deployment of this code: a multi-source leak in
two response handlers, and unsanitized shell interpolation of a source-IP value in two others.
Neither reached this paper's reported evaluation data (verified directly against the harness's own
trial-isolation design and the fact that every source-IP value the harness passes is a registry- or
Mininet-derived address, never externally supplied text), and both are documented here alongside
what would need fixing before any real deployment, rather than treated as already resolved.

\section{Conclusion}

We presented a registry-driven residential IoT defense system comparing three response-selection
policies against naive baselines and an external IDS on real, captured Mininet traffic, with a
statistical and ablation treatment -- including a formal paired significance test -- uncommon in
comparable systems papers. Our central positive result is that all three policies substantially and
reliably outperform naive fixed-response and no-response baselines ($99.6\%$ verified-response rate
versus $29.4\%$ and $5.9\%$), that policy \emph{choice} matters specifically at decision boundaries
rather than on confidently-classified traffic (McNemar's exact test, $p<1.2\times10^{-159}$ against
rule-based, $p=1.0$ between Stackelberg and PPO), and that this $99.6\%$ figure itself is the
product of finding and fixing six real, independently-verified defects, not an assumption. Our
central methodological result is that a scale-up is only worth running if it is treated as a test
the original finding can fail: a threefold sample-size increase first revealed our headline accuracy
had been optimistic, systematic investigation of every regressed condition then recovered and
substantially exceeded it, a source-rotation evasion claim did not survive a bug-fixed,
higher-powered re-test, and a training-budget fix that looked real on one isolated seed did not
generalize to the full population it was meant to help. We report all three corrections in full,
including the one that failed after already being committed, rather than retaining whichever
number was most favorable at each step. External validity beyond our single controlled environment
remains the primary open problem for future work; internal validity, in this paper, was earned
through repeated, disclosed self-correction rather than assumed from a single successful run.

\bibliographystyle{IEEEtran}
\bibliography{references}

\end{document}
